%(BEGIN_QUESTION)
% Copyright 2007, Tony R. Kuphaldt, released under the Creative Commons Attribution License (v 1.0)
% This means you may do almost anything with this work of mine, so long as you give me proper credit

Noen instrumentteknikere synes det er nyttig å tenke på integralinnstillingen ($\tau_i$) som et mål på regulatorens {\it utålmodighet}. Hva menes med denne analogien? I samme tankegang: er aggressiv integralvirkning mer ønskelig for regulering av prosesser som reagerer raskt av natur, eller prosesser som reagerer langsomt av natur?

\underbar{file i01611}
%(END_QUESTION)





%(BEGIN_ANSWER)

{\it Utålmodighet} er en treffende analogi for integralvirkning, fordi den knytter regulatorens utgang til hvor mye {\it tid} det finnes et avvik mellom PV og SP. Tilsvarende kan aggressiv integralvirkning egne seg for regulering av {\it hurtigresponderende} prosesser, men ikke for regulering av langsomt responderende prosesser.

%(END_ANSWER)





%(BEGIN_NOTES)

Dette betyr ikke at integralvirkning ikke kan brukes i langsomt responderende prosesser. Tvert imot er en fornuftig bruk av integralregulering den eneste måten å automatisk eliminere stasjonært avvik ved ren proporsjonalregulering. Poenget er bare at langsomme prosesser trenger {\it moderate} mengder integralvirkning, ikke for mye.

%INDEX% Control, integral: analogous to ``impatience''

%(END_NOTES)
